\section{Congruences}




\begin{definition}
Soit $n \ge 2$ un entier. On dit que $a$ est congru à $b$ modulo $n$,
si $n$ divise $b-a$. Dans ce cas on note :
$$a \equiv b \pmod n.$$

En d'autres termes : \eh{1}$a \equiv b \pmod n \quad \ssi \quad \exists k \in \Z \eh{2}\textrm{tel que}\eh{2} a = b +kn$.
\end{definition}


Remarquons que $n$ divise $a$ si et seulement si $a\equiv 0 \pmod n$.

\vspace{3mm}

\begin{exemples}{1mm}
	\item $59 \equiv 17 \pmod 7 \quad \text{car} \quad 59 = 7 \cdot 8 + \boxed{3} \quad \text{et} \quad 17 = 7 \cdot 2 + \boxed{3}$
	
	$59 \text{ est congru à } 17 \text{ modulo } 7 \text{ car ils ont le même reste.}$
	
	\item $3 \equiv -1 \pmod 2 \quad \text{car} \quad 3 = 2 \cdot 1 + \boxed{1} \quad \text{et} \quad -1 = 2 \cdot (-1) + \boxed{1}$
	
	\item$-14 \equiv -2 \pmod 4 \quad \text{car} \quad -14 = 4 \cdot (-4) + \boxed{2} \quad \text{et} \quad -2 = 4 \cdot (-1) + \boxed{2}$
	
	\item $-14 \not\equiv -2 \pmod 5 \quad \text{car} \quad -14 = 5 \cdot (-3) + \boxed{1} \quad \text{et} \quad -2 = 5 \cdot (-1) + \boxed{3}$
	
	$-14 \text{ n'est pas congru à } -2 \text{ modulo } 5 \text{ car ils n'ont pas le même reste.}$
	
	
	\item $-15 \not\equiv -2 \pmod 4 \quad \text{car} \quad -15 = 4 \cdot (-4) + \boxed{1} \quad \text{et} \quad -2 = 4 \cdot (-1) + \boxed{2}$
	
	$-15 \text{ n'est pas congru à } -2 \text{ modulo } 4 \text{ car ils n'ont pas le même reste.}$
\end{exemples}


\begin{proposition}
~

\vspace{-8mm}
\begin{enumerate}
  \item La relation \og congru modulo $n$ \fg{} est une relation d'équivalence :
  
  \vspace{-1mm}
  \begin{enumerate}[label=\roman*),left=0mm]
    \item (Réflexivité) $a \equiv a \pmod n$,
    \item (Symétrie) si $a \equiv b \pmod n$ alors $b \equiv a \pmod n$,
    \item (Transitivité) si $a \equiv b \pmod n$ et $b \equiv c \pmod n$ alors $a \equiv c \pmod n$.
  \end{enumerate}
  \item $a \equiv b \pmod n \iff a - b \equiv 0 \pmod n \iff n \mid (a - b) \iff a = nk + b \quad k \in \mathbb{Z}$
  
  \item Si $a \equiv b \pmod n$ et $c \equiv d \pmod n$ alors  la congruence est compatible avec :
  
  \begin{enumerate}[label=\roman*),left=0mm]
  	\item (l’addition) : $a \equiv a \pmod n$, $a+c \equiv b+d \pmod n$,
  	\item (la multiplication) : $a \equiv a \pmod n$, $a\cdot c \equiv b\cdot d \pmod n$,
  	\item (l'exponentiation) :  pour tout $k \geqslant 0$\; on a\; $a^k \equiv b^k \pmod n$.
  \end{enumerate}
  
\end{enumerate}
\end{proposition}



\vspace{4mm}
\noindent{\it Démonstration:}
\begin{enumerate}
  \item \textbf{1.i)} A voir que : $a \equiv a \pmod n \quad (\text{\textit{réflexive}})$
  
  Evident : Il existe des nombres $q \text{ et } r$ uniques tel que : $a = nq + r \quad \text{avec } 0 \le r < n$
  
  \vspace{0.3cm}
  
  \textbf{1.ii)} A voir que : Si $a \equiv b  \pmod n$ alors $b \equiv a  \pmod n \quad (\text{\textit{symétrique}})$
  
  Evident : Si $a \equiv b  \pmod n$ alors existe des nombres $q, q' \text{ et } r$ uniques tel que :
  
  $$a = nq + r \quad \text{et} \quad b = nq' + r \quad \text{avec } 0 \le r < n \quad \text{donc } b \equiv a  \pmod n$$
  
  \vspace{0.3cm}
  
  \textbf{1.iii)} A voir que : Si $a \equiv b  \pmod n$ et $b \equiv c  \pmod n$ alors $a \equiv c  \pmod n \quad (\text{\textit{transitive}})$
  
  $$\text{Si } a \equiv b  \pmod n \implies \begin{cases} a = nq + r & 0 \le r < n \\ b = nq' + r & 0 \le r < n \end{cases} \quad r \text{ unique}$$
  
  $$\text{et } b \equiv c  \pmod n \implies \begin{cases} b = nq' + r & 0 \le r < n \\ c = nq'' + r & 0 \le r < n \end{cases} \quad r \text{ unique}$$
  
  On constate que $a$ et $c$ ont le même reste lorsque l'on effectue une division euclidienne par $n$ d'où $a \equiv c  \pmod n$.
  \item  En Utilisant la définition,
  \begin{align*}
  	a \equiv b  \pmod n &\iff \begin{cases} a = nq + r \\ b = nq' + r \end{cases} \quad q \text{ et } q' \in \mathbb{Z} \text{ et } 0 < r \le n \\
  	&\iff a - b = (nq + r) - (nq' + r)\\
  	&\iff a - b = n \underbrace{(q - q')}_{\in \mathbb{Z}} + 0\\
  	&\iff a - b = n \underbrace{(q - q')}_{\in \mathbb{Z}} + 0\\
  	&\iff n \mid (a - b)\\
  	&\iff a - b \equiv 0  \pmod n\\
  	&\iff a - b = nk + 0 \quad k \in \mathbb{Z}\\
  	&\iff a = nk + b \quad k \in \mathbb{Z}
  \end{align*}

  \item i) Par hypothèse :
  
  $$a \equiv b  \pmod n \quad \text{et} \quad c \equiv d  \pmod n \overset{P2)}{\iff} \exists\, k, k' \in \mathbb{Z} \text{ tel que : } a = kn + b \quad \text{et} \quad c = k'n + d .$$
  
  \indent En additionnant ces deux égalités, on obtient :  
  $a + c = \underbrace{(k + k')}_{\in \mathbb{Z}} n + b + d \iff a + c \equiv b + d  \pmod n$\\
  
  \smallskip
  
  ii) Soit $a$ et $b$ tels que $a \equiv b \pmod n$, donc il existe $k \in \Z$ tel que $a=b + kn$.
  
%  \smallskip
   \indent Soit $c$ et $d$ tels que $c \equiv d \pmod n$, donc il existe $\ell \in \Z$ tel que $c \equiv d + \ell n$.
   
 \indent \indent Alors $a\fois c = (b+kn)\fois(d+\ell n) = bd + (b \ell + d k + k\ell n)\fois n$ qui est bien de la forme
$bd + m n$ avec $m\in \Z$. On conclut donc que  $a\fois c \equiv b\fois d \pmod n$.

   \smallskip
    iii) C'est une conséquence du point précédent : avec $a=c$ et $b=d$ on obtient
$a^2 \equiv b^2 \pmod n$. On continue par récurrence.
\end{enumerate}

\vspace{5mm}

La compatibilité de la relation de congruence avec l'addition, la multiplication et l’exponentiation permet de grandement simplifier les calculs de reste.

\vspace{5mm}

\begin{exemples}{1mm}
	\item  On a $50^{100} \equiv 1 \pmod 7$.
	 
	 En effet, comme  $50 \equiv 1 \pmod 7$ (car $50 = 7  \cdot 7 + 1$).	\\
	 De plus, d'après la compatibilité avec les puissances, on a : $50^{100} \equiv 1^{100} \equiv 1 \pmod 7$.
	
	
	\vspace{0.3cm}
	
	\item  On a $100^{3} \equiv 1 \pmod 7$.
	
	 En effet, comme $100 \equiv 2 \pmod 7$ (car $100 \equiv 2\cdot 50 \equiv 2\cdot 1 \pmod 7$), d'après la compatibilité avec les puissances, on a :
	$100^3 \equiv 2^3 \equiv 8 \equiv 1 \pmod 7$.
	
	\vspace{0.3cm}
	
	\item  On a $50^{100} + 100^{100} \equiv  3 \pmod 7$.
	
	 En effet, comme $100^{100} = 100^{3  \cdot 33 + 1} = \left(100^3\right)^{33}  \cdot 100$, donc d'après la compatibilité avec les puissances et la multiplication, on a : $100^{100} \equiv 1^{33} \cdot 2 \equiv 2 \pmod 7$.
	
	D'après la compatibilité avec l'addition, on a alors : $50^{100} + 100^{100} \equiv 1 + 2 \equiv 3 \pmod 7$.
	
	Le reste est 3.
\end{exemples}

\vspace{5mm}

De même, on peut utiliser ces relations pour déterminer si on nombre est un diviseur (ou un multiple) d'un autre.

\vspace{5mm}

\begin{exemple}
	On peut montrer que $\forall n \in \mathbb{N}$, le nombre $3^{n+3}-4^{4 n+2}$  est divisible par 11.
	
	\smallskip \smallskip
	
	
	En effet, on a : $3^{n+3} = 3^n \cdot 3^3 = 27 \cdot 3^n$.
	
	Or $27 \equiv 5 \pmod{11}$, donc d'après la compatibilité avec la multiplication, on a :
	
	$$\forall n \in \mathbb{N}, \quad 3^{n+3} \equiv 5 \cdot 3^n \pmod{11}$$
	
	On a : $4^{4n+2} = \left(4^4\right)^n \cdot 4^2$, or $4^2 \equiv 5 \pmod{11}$ donc $4^4 \equiv 5^2 \equiv 3 \pmod{11}$, donc :
	
	$$\forall n \in \mathbb{N}, \quad 4^{4n+2} \equiv 3^n \cdot 5 \pmod{11}$$
	
	On en déduit donc que :
	
	$$3^{n+3} - 4^{4n+2} \equiv 0 \pmod{11}$$
	
	La proposition est donc vérifiée pour tout entier naturel $n$.
\end{exemple}

%\begin{remarque}
%Pour trouver un \og bon \fg{} représentant de $a \pmod n$ on peut aussi faire la division euclidienne
%de $a$ par $n$ : $a = bn + r$ alors $a \equiv r \pmod n$ et $0 \le r < n$.
%\end{remarque}



\begin{exstheorie}
	\exo
	Trouver le plus petit nombre supérieur ou égal à 0 qui est congru à $a$ modulo $m$.
	
	
	\begin{enumerate}[label=\roman*),left=0mm]
		\item $a=3412$ et $m=4$
		\item $a=177$ et $m=9$
		\item $a=31$ et $m=31$
		\item $a=31$ et $m=25$
		\item $a=365$ et $m=5$
		\item $a=-3122$ et $m=3$
		\item $a=31458687312$ et $m=10$
		\item $a=-259874629$ et $m=10$
	\end{enumerate}
	
	
	
	\eexo
	\pagebreak[3]
	
	\exo
	\noindent Les affirmations suivantes sont-elles vraies ou fausses ? Justifier vos réponses.
	
	\begin{enumerate}
		\item[\textbf{a)}] $471 \equiv 11 \pmod {23}$
		\item[\textbf{b)}] $370 \equiv 17 \pmod {13}$
		\item[\textbf{c)}] $29 \equiv -121 \pmod 5$
		\item[\textbf{d)}] $-30 \equiv -1 \pmod {26}$
		\item[\textbf{e)}] L'égalité $914 = 19 \cdot 47 + 21$ permet de déduire les congruences :
		
		$914 \equiv 21 \pmod {19} \quad$ et $\quad 914 \equiv 21 \pmod {47}$
		
		\item[\textbf{f)}] Soit $n \in \mathbb{N}$ avec $n \ge 2$ et $a \in \mathbb{N}^*$. \\
		Si $a \mid b \iff b \equiv 0 \pmod a$
		\item[\textbf{g)}] Soit $n \in \mathbb{N}$ avec $n \ge 2$ et $a \in \mathbb{N}^*$. \\
		Si $b \equiv 0 \pmod a$ et $c \equiv 0 \pmod a$ alors $\alpha \cdot b + \beta \cdot c \equiv 0 \pmod a$ avec $\alpha, \beta \in \mathbb{N}^*$ \\
		Autrement dit : Si $a$ divise $b$ et $a$ divise $c$ alors $a$ divise toute combinaison linéaire de $b$ et de $c$
	\end{enumerate}
	\eexo
	\pagebreak[3]	
	\exo
	Trouver les restes de la division euclidienne par $7$ des nombres :
	$351^{12} \cdot 85^{15}$ et $\quad 16^{12} - 23^{12}$.
	
	
	
	\eexo
	\pagebreak[3]
	
	
	\exo
	Démontrer que pour tout entier naturel $n$, $5^{2n} - 14^n$ est divisible par $11$.
	
	\eexo
	\pagebreak[3]
\end{exstheorie}

\begin{exstheorie}

\exo
Trouver tous les nombres compris entre $1950$ et $2000$ qui sont congrus à $a$ modulo $m$.

\begin{enumerate}[label=\roman*),left=0mm]
\item $a=1$ et $m=13$
\item $a=1776$ et $m=40$
\item $a=1929$ et $m=15$
\end{enumerate}


\eexo
\pagebreak[3]



\exo
Démontrer par récurrence que $6\cdot 4^n \equiv 6 \pmod{9}$ pour tout $n\geqslant 0$.



\eexo
\pagebreak[3]



\exo
Trouver le reste de la division par $13$ du nombre $100^{1000}$.


\eexo
\pagebreak[3]


\exo
Démontrer que le nombre $7^n+1$ est divisible par $8$ si $n$ est impair ; dans le cas où $n$ est pair, donner le reste de la division.
\end{exstheorie}


\medskip

%\begin{exemple}{0.9} \textbs{Exponentiation rapide (\guillemets{square-and-multiply})}
%
%\medskip
%Les calculs bien menés avec les congruences sont souvent très rapides. L'exponentiation rapide est un algorithme Par exemple
%on souhaite calculer $2^{21} \pmod {37}$ (plus exactement on souhaite trouver $0 \le r < 37$ tel que
%$2^{21} \equiv r \pmod {37}$).
%Plusieurs méthodes :
%\begin{enumerate}
%\item On calcule $2^{21}$, puis on fait la division euclidienne de $2^{21}$ par $37$, le reste est notre résultat.
%C'est laborieux !
%\item On calcule successivement les $2^k$ modulo $37$ :
%$2^1 \equiv 2 \pmod {37}$, $2^2 \equiv 4 \pmod {37}$, $2^3 \equiv 8 \pmod {37}$, $2^4 \equiv 16 \pmod {37}$,
%$2^5 \equiv 32 \pmod {37}$. Ensuite on n'oublie pas d'utiliser les congruences :
%$2^6 \equiv 64 \equiv 27 \pmod {37}$.
%$2^7 \equiv 2 \cdot 2^6 \equiv 2 \cdot 27 \equiv 54 \equiv 17 \pmod {37}$
%et ainsi de suite en utilisant le calcul précédent à chaque étape.
%C'est assez efficace et on peut raffiner : par exemple on trouve $2^{8} \equiv 34 \pmod {37}$
%mais donc aussi $2^{8} \equiv -3 \pmod {37}$ et donc
%$2^9 \equiv 2 \cdot 2^8 \equiv 2 \cdot(-3) \equiv -6 \equiv 31 \pmod{37}$,...
%
%\item Il existe une méthode encore plus efficace, on écrit l'exposant $21$ en base $2$ :
%$21 = 2^{4} + 2^{2} + 2^{0} = 16 + 4 + 1$. Alors
%$2^{21} =  2^{16} \cdot 2^{4} \cdot 2^{1}$.
%Et il est facile de calculer successivement chacun de ces termes car les exposants sont des puissances de $2$.
%Ainsi $2^8 \equiv  (2^4)^2 \equiv 16^2 \equiv 256 \equiv 34 \equiv -3 \pmod {37}$ et
%$2^{16} \equiv  \left(2^{8}\right)^2 \equiv (-3)^2 \equiv 9 \pmod{37}$.
%Nous obtenons $2^{21} \equiv  2^{16} \cdot 2^{4} \cdot 2^{1} \equiv 9 \times 16 \times 2 \equiv 288 \equiv 29 \pmod {37}$.
%\end{enumerate}
%\end{exemple}


\begin{definition}
Une classe d'équivalence \guillemets{$a$ modulo $n$} ($n\geqslant 2$), notée $[a]_n$, est l'ensemble de tous les nombres entiers dont la division euclidienne par $n$ donne le même reste que la division de $a$ par $n$.

\smallskip
On note la classe d'équivalence de $a$ \textit{modulo} $n$ : $[a]_n$
\end{definition}


\medskip
\begin{exemple}
$[0]_3=[3]_3=[6]_3=\ldots$

\medskip
$[1]_3=[4]_3=[7]_3=\ldots$

\medskip
$[2]_3=[5]_3=[8]_3=\ldots$
\end{exemple}


%\begin{definition}
%Un groupe est la donnée d'une ensemble $G$ muni d'une opération $\ast$, qu'on note $(G,\ast)$ et qui possède les propriétés suivantes:
%\begin{enumerate}[label=\textbf{(\arabic*)},left=0mm,itemsep=1mm]
%\item \textbf{L'opération est interne :} Si $g,h \in G$ alors $g\ast h\in G$
%\item \textbf{Existence d'un élément neutre :} $\exists\eh{1} e\in G$ tel que $g\ast e=e\ast g=g \eh{2}\forall\eh{1} g\in G$.
%\item \textbf{Existence de l'inverse :} $\forall\eh{1} g\in G \virg[1] \exists\eh{1} \overline{g}\in G$ tel que $g\ast \overline{g}=\overline{g}\ast g=e$.
%\item \textbf{Associativité :} $\forall\eh{1} g, h, j$ on a $(g\ast h)\ast j=g\ast (h\ast j)$.
%\end{enumerate}
%
%\end{definition}

\bigskip
\begin{proposition}[\normalsize Équations de congruences]
Soit $a \in \Z^*$, $b \in \Z$ fixés et $n \ge 2$.
Considérons l'équation
{$ax \equiv b \pmod n$} d'inconnue $x \in \Z$.
\begin{enumerate}[label=\textbf{(\arabic*)},left=0mm,itemsep=1mm]
  \item Il existe des solutions si et seulement si $\pgdc(a,n) | b$.
  \item Les solutions sont de la forme $x = x_0 + \ell \frac{n}{\pgdc(a,n)}$, où $\ell \in \Z$
et $x_0$ est une solution particulière.

Il existe donc $\pgdc(a,n)$ classes d'équivalences solutions.
\end{enumerate}
\end{proposition}

\vspace{1mm}
\textit{Démonstration :}


\begin{enumerate}[label=\textbf{(\arabic*)},left=0mm,itemsep=1mm]
%\item $x \in \Z \text{ est solution de l'équation }   ax \equiv b \pmod{n}$
%$\begin{aligned}[t]
% \iff\   & \exists k \in \Z \quad ax = b +kn \\
% \iff\   & \exists k \in \Z \quad ax - kn = b \\
% \iff\   &  \pgdc(a,n) | b \quad \text{(proposition \ref*{prop:equations_diophantiennes})}\\
%\end{aligned}$
\item Il suffit de remarquer que nous avons les équivalences suivantes :

\smallskip
$x \in \Z \text{ est solution de l'équation }   ax \equiv b \pmod{n}$
$\setlength{\jot}{2mm}\begin{aligned}[t]
 \iff\   & \exists k \in \Z \eh{1}\text{ tel que }\eh{1}  ax = b +kn \\
 \iff\   & \exists k \in \Z \eh{1}\text{ tel que }\eh{1} ax - kn = b \\
 \iff\   &  \pgdc(a,n) | b \quad \text{(proposition \ref*{prop:equations_diophantiennes})}\\
\end{aligned}$
\smallskip
%
%Nous avons juste transformé notre équation $ax\equiv b \pmod n$
%en une équation $ax-kn=b$ étudiée auparavant (voir section ),
%seules les notations changent :
%$\textcolor{magenta}{a}\textcolor{red}{u}+\textcolor{blue}{b}\textcolor{orange}{v}=\textcolor{green}{c}$ devient
%$\textcolor{magenta}{a}\textcolor{red}{x}-\textcolor{orange}{k}\textcolor{blue}{n}=\textcolor{green}{b}$.

\item Supposons qu'il existe des solutions. On note  $d = \pgdc(a,n)$ et on peut écrire $a=da'$, $n = dn'$
et $b = db'$ (car, par le premier point, $d|b$). Remarquons également que $\pgdc(a';n')=1$.

\smallskip
L'équation de congruence {$ax \equiv b \pmod n$} d'inconnue $x$ se ramène à l'équation diophantienne $ax-kn=b$ d'inconnues $x,k \in \Z$ qui est équivalente à
l'équation $$a'x-kn'=b'$$ 
Comme $\pgdc(a';n')=1 | b$, alors l'équation possède des solutions et sa résolution (voir les équations diophantiennes) donne que
$$x = x_0 + \ell n'$$
où $x_0$ est une solution particulière et $\ell$ parcourt $\Z$ (les $k$ ne nous intéressent pas ici).

\smallskip
Ainsi les solutions $x\in \Z$ sont de la forme 
$$x = x_0 + \ell \frac{n}{\pgdc(a,n)} \virg[1]\ell \in \Z$$
et modulo $n$ cela donne bien $\pgdc(a,n)$ classes d'équivalences distinctes.\hfill$\blacksquare$
\end{enumerate}




\begin{exemplesecable}
Résolvons l'équation $9x \equiv 6 \pmod{24}$.

\medskip
Comme $\pgdc(9,24)=3$ divise $6$, la proposition ci-dessus nous affirme qu'il existe des solutions. Il faudra toujours refaire  les calculs et non reprendre directement la formule.

\medskip
Trouver $x$ tel que $9x \equiv 6 \pmod{24}$ est équivalent à trouver $x$ et $k$ tels que
$9x = 6 + 24k$. Cette dernière équation peut s'écrire sous la forme $9x-24k=6$. (qui possède bien des solutions car $\pgdc(9,24)=3$ divise $6$)

\medskip
En divisant par le pgdc on obtient l'équation diophantienne équivalente :
 $$3x-8k=2.$$

Pour le calcul du pgdc et d'une solution particulière nous utilisons normalement l'algorithme d'Euclide et sa remontée.
Ici il est facile de trouver une solution particulière :
$$(x_0=6\pvirg{0.5}k_0=2)$$

\medskip
On termine comme pour les équations diophantiennes. Si $(x\pvirg{0.5}k)$ est une solution
de l'équation alors par soustraction on obtient $3(x-x_0)-8(k-k_0)=0$. En suivant le raisonnement de résolution d'une équation diophantienne (laissée en exercice ici), on trouve :

\vspace{-3mm}
$$x = x_0 + 8\ell \eh{1}\text{ avec }\eh{1} \ell \in \Z$$ 

(le terme $k$ ne nous intéresse pas).
Nous avons donc trouvé les $x$ qui sont solutions de $9x-24k=6$, ce qui équivaut à $9x \equiv 6 \pmod{24}$. 

\medskip
Les solutions sont donc de la forme $x=6+ 8\ell$, ce qui en classes d'équivalences modulo $24$ donne :

$$S=\left\lbrace [6]_{24} \pvirg{1} [14]_{24} \pvirg{1} [22]_{24} \right\rbrace$$
\end{exemplesecable}





\begin{exstheorie}
\exo
Trouver les restes de la division euclidienne par 11 des nombres suivants :
\begin{tasks}[label=\alph*),item-indent=6mm,column-sep=0mm ,after-item-skip=5mm,label-offset=2mm,label-align=right](3)
\task $12^{15}$
\task $10^7$
\task $78^{15}$
\task $13^{12}$
\task $(-2)^{19}$
\end{tasks}





\eexo
\pagebreak[3]





\exo
~

\vspace{-6mm}
\begin{tasks}[label=\alph*),item-indent=6mm,column-sep=0mm ,after-item-skip=3mm,label-offset=2mm,label-align=right](1)
\task Vérifier que $2^4 \equiv -1 \pmod{17}$ \eh{2}et que \eh{2}\\ $6^2 \equiv 2 \pmod{17}$.
\task Déterminer la classe d'équivalence de $\left[1532\right]_{17}$ \et{2} $\left[346\right]_{17}$.
\task En déduire le reste de la division par 17 des nombres $1532^{20}$ \et{2} $346^{12}$.
\end{tasks}





\eexo
\pagebreak[3]





\exo
Résoudre les équations de congruences suivantes :
\begin{tasks}[label=\alph*),item-indent=6mm,column-sep=0mm ,after-item-skip=3mm,label-offset=2mm,label-align=right](1)
\task $3x \equiv 4 \pmod{7}$
\task $4x \equiv 14 \pmod{30}$
\end{tasks}



\columnbreak
%\eexo




\exo
~

\vspace{-7mm}
\begin{tasks}[label=\alph*),item-indent=6mm,column-sep=0mm ,after-item-skip=5mm,label-offset=2mm,label-align=right](1)
\task Montrer que le reste de la division euclidienne par $8$ du carré de tout
nombre impair est $1$.
\task Montrer de même que tout nombre pair vérifie :
$$x^{2}\equiv 0\pmod{8} \ou{4} x^{2}\equiv 4\pmod{8}$$\vspace{-4mm}
\task Soient $a$, $b$, $c$ trois entiers impairs.\\
Déterminer le reste modulo $8$ de :
\begin{enumerate}[itemsep=1mm]
\item[] $\bm{a^{2}+b^{2}+c^{2}}$
\item[] $\bm{2(ab+bc+ac)}$
\end{enumerate}
{\footnotesize \textsl{\textbs{Indications :} Utiliser l'item a) pour le premier et écrire $a$, $b$ et $c$ comme des nombres impairs pour le deuxième.}} 
\task En déduire que :
$$a^{2}+b^{2}+c^{2} \et{4} 2(ab+bc+ac)$$
ne sont pas des carrés.

\medskip
Conclure que $ab+bc+ac$ non plus.
\end{tasks}
\end{exstheorie}

